\subsection{Автоматизация проверки}

Для запуска верификации используется скрипт \code{run.py} на языке Python.
Он работает со стандартной библиотекой и объединяет подготовку файлов,
вызов SPIN, компиляцию верификаторов и обработку результатов. Проверка
запускается командой \code{python run.py} из корня проекта.

\subsubsection{Подготовка файлов и запуск инструментов}

Работа начинается в функции \code{main}. В каталоге
\code{verification_results} создается отдельная папка \code{run_*},
чтобы артефакты нового запуска не перезаписывали предыдущие.
Функция \code{prepare_files} копирует
в него модель, свойства и конфигурацию, после чего читает пути к SPIN
и GCC из \code{config.json}. Имена проверяемых свойств извлекаются
из объявлений \code{ltl} регулярным выражением.

\listing[Python]{prepare_files.py}{Подготовка входных данных проверки}

Функция \code{run} выполняет внешнюю команду через \code{subprocess.run},
сохраняет ее полный вывод в отдельный лог и проверяет код завершения.
Для каждой команды установлен лимит в 180 секунд. Ошибка запуска,
ненулевой код или превышение времени не позволяют считать проверку
успешной.

Далее \code{build_verifiers} сохраняет версии инструментов, вызывает
\code{spin -a properties.pml} и дважды компилирует полученный \code{pan.c}.
Один исполняемый файл предназначен для структурной проверки модели,
второй используется для проверки LTL-свойств. Основные команды приведены ниже.

\listing[]{commands.txt}{Генерация и запуск верификаторов}

Структурная проверка выявляет ситуации, в которых ни один процесс
не может выполнить следующий шаг, хотя работа системы не завершена.
Например, Master может ожидать уведомление, которое ни один Slave
уже не способен отправить. SPIN рассматривает такую остановку
как недопустимое конечное состояние. Эта проверка нужна, чтобы отдельно
выявить ошибки обмена сообщениями перед проверкой LTL-свойств.

При компиляции верификаторов используются следующие параметры~\cite{pan}:
\begin{itemize}
    \item \code{-DNOCLAIM} исключает автомат \code{never claim}
    из структурной проверки. Поэтому она выполняется независимо
    от выбранной LTL-формулы.
    \item \code{-DSAFETY} отключает поиск циклов и сокращает затраты
    на проверку. Для обнаружения недопустимого конечного состояния
    достаточно конечной последовательности шагов, которая к нему приводит.
    \item \code{-DCOLLAPSE} сжимает хранимые состояния без потери данных
    и уменьшает расход памяти. Проверка остается полной.
    Приближенный режим \code{BITSTATE}, который может пропустить
    часть состояний, не используется.
\end{itemize}

При запуске LTL-верификатора параметр \code{-N} выбирает формулу по имени,
а \code{-a} включает поиск принимающих циклов. Такой цикл
описывает бесконечное выполнение модели, нарушающее выбранную формулу,
например бесконечное ожидание вместо завершения транзакции.

Максимальная глубина обхода задана как 1\,000\,000 шагов. Параметр
\code{-b} делает превышение глубины ошибкой, а \code{-q} в структурной
проверке требует пустых каналов в допустимом конечном состоянии.
\code{-w22} задает начальный размер хеш-таблицы поиска.

\subsubsection{Обработка результатов и контрпримеров}

Сам по себе нулевой код завершения программы еще не является
подтверждением свойства. Функция \code{check_search_output} читает число
ошибок, проверяет отсутствие сообщений о нехватке ресурсов и убеждается,
что верификатор действительно выполнял полный поиск в нужном режиме.

\listing[Python]{check_search_output.py}{Проверка достоверности результата поиска}

Для структурной проверки должен быть включен поиск недопустимых конечных
состояний. Для LTL должны быть активны \code{never claim} и поиск
принимающих циклов. Затем функция \code{verify} различает найденное
нарушение, успешную проверку и техническую ошибку.

\listing[Python]{verify.py}{Формирование статуса одной проверки}

Если обнаружено нарушение, \code{save_counterexample} сохраняет
\code{properties.pml.trail} под именем свойства и запускает SPIN
с параметрами \code{-t -p -g -l -k}. В отдельный лог записывается
пошаговое воспроизведение с состояниями процессов и переменных.
Если поиск не завершился или возникла ошибка инструмента,
результат получает статус \code{ERROR}, а не \code{PASS}.

Функция \code{run_checks} сначала проверяет структуру модели, затем
последовательно запускает все найденные LTL-свойства. В конце сохраняется
\code{report.txt} с результатами и общим количеством \code{PASS},
\code{FAIL} и \code{ERROR}. Код завершения скрипта равен 0 при полном
успехе, 1 при нарушении свойства и 2 при технической ошибке. При наличии
и нарушения, и технической ошибки возвращается 2.
